\begin{abstract}
A web agent can read a page as text, look at a screenshot, or do both, and deployments usually fix one choice for every task. We measure that choice across three sites of VisualWebArena and WebArena with four vision-language backbones, using text-only variants as controls and rerunning identical conditions to see which differences survive.
% sources: routing_three_arm, visual_intent_routing, task_mode_interaction_null, routing_gain_upper_bounds
Three findings separate questions that are usually merged. First, no observation choice is best everywhere, and the screenshot has a real, reproducible value that is predictable in advance on one site: a zero-cost rule over the task text picks out classifieds tasks on which adding the screenshot alone raises success by 25--30 points for the two strongest backbones (11 for a 4B model), and the gap reappears on a full rerun. Second, which representation suits which task is a reproducible property of the task, beyond an additive model of task difficulty, yet a single run measures it with reliability of only 0.31--0.50. Third, two lightweight routers that choose a representation per task from pre-run features, with their operating point fixed on training data, show no reliable gain over the fixed choices and their random mixtures. We describe the rerun-controlled protocol that keeps these three quantities apart.
\end{abstract}
